% 全文参考文献；从问题一冲突消解章节移出，以便在各问正文后统一排版。
\begin{thebibliography}{99}
\bibitem{diakoulaki1995} DIAKOULAKI D, MAVROTAS G, PAPAYANNAKIS L. Determining objective weights in multiple criteria problems: The CRITIC method[J]. Computers \& Operations Research, 1995, 22(7): 763--770. DOI: \href{https://doi.org/10.1016/0305-0548(94)00059-H}{10.1016/0305-0548(94)00059-H}.
\bibitem{giarlotta2001} GIARLOTTA A. Multicriteria compensability analysis[J]. European Journal of Operational Research, 2001, 133(1): 190--209. DOI: \href{https://doi.org/10.1016/S0377-2217(00)00198-3}{10.1016/S0377-2217(00)00198-3}.
\bibitem{yu2016} 余高锋, 李登峰, 邱锦明, 等. 考虑偏好冲突的多类评价信息群体决策方法[J]. 控制与决策, 2016, 31(11): 2013--2018. DOI: \href{https://doi.org/10.13195/j.kzyjc.2015.0948}{10.13195/j.kzyjc.2015.0948}.
\bibitem{wang2023} 王雷, 王晓玲, 张君, 等. 考虑融合权重优化与冲突信息来源的大坝安全综合评价方法[J]. 清华大学学报(自然科学版), 2023, 63(10): 1566--1575. DOI: \href{https://doi.org/10.16511/j.cnki.qhdxxb.2023.22.039}{10.16511/j.cnki.qhdxxb.2023.22.039}.
\bibitem{chen2026} 陈雪龙, 王子睿, 李苗苗. 多周期应急资源调配主体冲突消解方法研究[J]. 中国管理科学, 2026, 34(3): 226--241. DOI: \href{https://doi.org/10.16381/j.cnki.issn1003-207x.2023.2186}{10.16381/j.cnki.issn1003-207x.2023.2186}.
\bibitem{penedo2024} PENEDO G, KYDLÍČEK H, BEN ALLAL L, et al. The FineWeb datasets: Decanting the web for the finest text data at scale[EB/OL]. 2024. \url{https://arxiv.org/abs/2406.17557}.
\bibitem{li2024} LI J, FANG A, SMYRNIS G, et al. DataComp-LM: In search of the next generation of training sets for language models[C]//Advances in Neural Information Processing Systems. 2024. \url{https://arxiv.org/abs/2406.11794}.
\end{thebibliography}
